\documentclass[11pt]{article}
\usepackage[margin=1in]{geometry}
\usepackage{amsmath,amssymb,amsthm,booktabs,array}
\usepackage[T1]{fontenc}
\usepackage{lmodern}
\usepackage{microtype}
\usepackage{hyperref}
\hypersetup{colorlinks=true,linkcolor=blue,urlcolor=blue}
\setlength{\emergencystretch}{3em}
\newtheorem{theorem}{Theorem}
\newtheorem{lemma}{Lemma}
\newcommand{\cI}{\mathcal I}
\newcommand{\norm}[1]{\left\lVert #1\right\rVert}
\newcommand{\XiN}{\mathsf{Xi}}
\title{Cross-Scale Endpoint-Normalized Radius in a Finite V59 Residual}
\author{Liang Wang\\School of Mathematics and Statistics\\Huazhong University of Science and Technology (HUST)\\Wuhan, China}
\date{27 August 2026}
\begin{document}
\maketitle
\begin{abstract}
TPC-269 showed that a registered growing-cutoff proxy and a convex profile path
still cross the quarter threshold, but it did not measure the size of the
residual radius. We introduce the dimensionless finite observable
\(\XiN_{N,\theta}=(R_{N,\theta}^{2})^3/N^{10}\), equivalently
\((R_{N,\theta}/N^{5/3})^6\), so that the endpoint-scale normalization is
rational once the upstream interval for \(R^2\) is certified. On the same
source-compatible finite registry, four dyadic ratio intervals exhibit the
separated pattern \(\text{DROP--RISE--RISE--DROP}\): the ratio for
\(96\!\to\!192\) is above \(23\), while that for \(64\!\to\!128\) is below
\(1/4\). Three matched profile controls remain in \((1/2,3/4)\). This is a
finite normalization audit and a scoped stability obstruction; it is not an
asymptotic radius estimate, a power-saving theorem, or a twin-prime result.
\end{abstract}

\section{Position and claim firewall}
The TPC-266 typed compiler leaves the Schur residual radius as a separate lane.
TPC-267 instantiated the literal prime-shell object at finite scale, and
TPC-268 showed that declared finite cutoff changes can alter its quarter-sector
classification. TPC-269 then used \(z_N=\lfloor\log N\rfloor\) and an affine
two-profile path, finding a finite profile flip. The remaining minimal question
is a scale question: does the radius itself have a stable normalization on the
same finite interface?

The status of this paper is
\begin{center}
\texttt{NUMERICALLY\_CERTIFIED\_FINITE\_CROSS\_SCALE\_RADIUS\_NORMALIZATION\_AUDIT}.
\end{center}
The word ``finite'' is part of the claim. The registered scales are a diagnostic
registry, not an asymptotic sequence with a proved source-level uniformity law.
In particular, no observed ratio below or above one is promoted to an eventual
lower or upper bound.

\section{The frozen finite residual}
Let \(N\) be one of the registered even scales and write
\[
 \cI_N=\{N/2+1,\ldots,N\}, \qquad
 \mathcal Q_Q=\{q:q\ \mathrm{prime},\ Q<q\leq 2Q\}.
\]
The source and shifted-prime comparison are the finite objects used in TPC-269:
\[
 \beta_N(t)=\frac{\Lambda(t)}{\log t}
 -\sum_{d\mid t,\ d^{400}\leq N^{133}}\mu(d),
 \qquad w_N(u)=\Lambda(u+2)-b_N^{(z_N)}(u),
\]
where \(z_N=\lfloor\log N\rfloor\) on the registry and
\[
 b_N^{(z)}(u)=C_2^{(z)}1_{2\nmid u}
 \prod_{p\leq z}\frac p{p-1}
 \prod_{\substack{p\mid u\\p>z}}\frac{p-1}{p-2}.
\]
As before, the value is zero when \(u+2\) has a prime factor at most \(z\),
and the Euler product is enclosed by the finite tail protocol.

For \(s=1,2\), let
\[
 K_{H,s}(h)=\left(1+\frac{h^2}{H^2}\right)^{-s}
\]
and retain the prime shell, unit masks, deleted diagonal, and centered residue
factor in
\[
 A_s(u,t)=1_{u\ne t}\sum_{q\in\mathcal Q_Q}qK_{H,s}(u-t)1_{q\nmid ut}
 \left(1_{u\equiv t\pmod q}-\frac1{q-1}\right).
\]
For the base rows we use \(\theta=0\), so \(g_{N,0}=A_1\beta_N\). At three
matched scales we also use the source-compatible finite mixture
\[
 A_\theta=(1-\theta)A_1+\theta A_2, \qquad
 g_{N,\theta}=A_\theta\beta_N, \qquad \theta=1/2.
\]

Partition \(\cI_N\) into four equal blocks and let \(P_3\) be the orthogonal
projection onto the three contrasts
\[
 (1,1,-1,-1),\qquad (1,-1,0,0),\qquad (0,0,1,-1).
\]
The residual product is
\[
 R_{N,\theta}^2=\norm{(I-P_3)w_N}^2
                  \norm{(I-P_3)g_{N,\theta}}^2.
\]
All non-logarithmic terms are rational on a finite row. The logarithms and the
Euler constant use outward intervals with tail cutoff \(P=50000\).

\section{Endpoint normalization}
The endpoint budget in TPC-265 uses the baseline exponent \(5/3\) for a radius
lane. Directly taking \(N^{5/3}\) would introduce an unnecessary algebraic
number into a finite certificate. We instead define
\[
 \XiN_{N,\theta}:=\frac{(R_{N,\theta}^2)^3}{N^{10}}.
 \tag{1}
\]

\begin{lemma}[exact normalization identity]
If \(R_{N,\theta}^2>0\), then
\[
 \XiN_{N,\theta}=\left(\frac{R_{N,\theta}}{N^{5/3}}\right)^6.
\]
If \([r_-,r_+]\) is a positive outward interval for \(R^2\), then
\[
 \XiN\in\left[\frac{r_-^3}{N^{10}},\frac{r_+^3}{N^{10}}\right].
 \tag{2}
\]
\end{lemma}
\begin{proof}
The first identity follows from \(R^6=(R^2)^3\) and
\(N^{10}=(N^{5/3})^6\), and the second
follows because cubing is increasing on the positive half-line. Every endpoint
in (2) is rational in the finite interval engine.
\end{proof}

For positive intervals \([a_-,a_+]\) and \([b_-,b_+]\), the derived scale ratio
uses
\[
 \frac{[a_-,a_+]}{[b_-,b_+]}=
 \left[\frac{a_-}{b_+},\frac{a_+}{b_-}\right].
 \tag{3}
\]
We study \(D(a,b)=\XiN_{b,0}/\XiN_{a,0}\) for the four registered dyadic
pairs, and compare \(\XiN_{N,1/2}/\XiN_{N,0}\) at three scales.

\section{Certified finite result}
The six base rows use
\[
 (N,H,Q)=(64,15,4),(96,20,5),(128,24,5),
 (192,32,6),(256,38,6),(384,50,7).
\]
The resulting \(\XiN_{N,0}\) intervals are shown in Table~\ref{tab:base}; the
stored radius-square intervals are inherited from the exact finite operator
replay and are not floating-point point estimates.

\begin{table}[t]
\centering
\caption{Base endpoint-normalized sixth-power intervals.}
\label{tab:base}
\small
\begin{tabular}{rcc}
\toprule
N & \(z_N\) & \(\XiN_{N,0}\) interval\\
\midrule
64  &4& [8.81278791526e-7,\ 8.81449971308e-7]\\
96  &4& [4.95949095187e-7,\ 4.96030749164e-7]\\
128 &4& [2.04279432566e-7,\ 2.04322254881e-7]\\
192 &5& [1.18847779301e-5,\ 1.18871727315e-5]\\
256 &5& [1.46532173424e-6,\ 1.46559968348e-6]\\
384 &5& [9.54437330125e-6,\ 9.54594504625e-6]\\
\bottomrule
\end{tabular}
\end{table}

\begin{theorem}[finite cross-scale normalization audit]
On the declared registry, the four dyadic ratios are
\begin{center}
\begin{tabular}{ccl}
\(a\to b\)& \(D(a,b)\) interval & finite classification\\
\midrule
64\(\to\)128 & [0.231753859227, 0.231847466257] & drop below \(1/4\)\\
96\(\to\)192 & [23.9597604587, 23.9685339622] & rise above 16\\
128\(\to\)256 & [7.17162080603, 7.17448479796] & rise above 7\\
192\(\to\)384 & [0.802913654645, 0.803207691586] & drop below 1\\
\end{tabular}
\end{center}
Thus the dyadic pattern is
\[
 \mathrm{DROP\!-\!RISE\!-\!RISE\!-\!DROP}.
\]
At the matched scales \(N=96,128,256\), the profile ratios obey
\[
\begin{array}{c|c}
N&\XiN_{N,1/2}/\XiN_{N,0}\\ \hline
96 &[0.632933594160,0.633142026033]\\
128 &[0.716167972521,0.716468259088]\\
256 &[0.600909229748,0.601137218046]
\end{array}
\]
and hence all lie strictly between \(1/2\) and \(3/4\).
\end{theorem}
\begin{proof}
The finite operator and projection are the frozen TPC-269 objects. The interval
engine first certifies positive \(R^2\), applies (2), and then applies (3) to
the listed pairs. Every displayed classification is separated from its stated
threshold. The independent checker recomputes the source, operator, projection,
and normalization using a separate floating-point implementation; the stress
audit checks the registry and all threshold inequalities.
\end{proof}

The five adjacent base ratios give the supplementary pattern
\(\mathrm{DROP, DROP, RISE, DROP, RISE}\), with the largest finite rise at
\(128\to192\) (above 58). The profile controls are deliberately reported as
ratios at the same scale: they quantify a finite profile effect without being
mistaken for a cross-scale theorem.

\section{Interpretation and limits}
Equation (1) is useful because it turns the radius question into a dimensionless
positive quantity and makes cross-scale comparisons auditable. The main finite
obstruction is not the absolute size of any one row; it is the coexistence of a
large normalized rise and a strict normalized drop under one declared cutoff
rule and one base profile family.

Four promotions are prohibited. First, the six scales do not establish an
asymptotic sequence or a uniform source-level estimate. Second, a finite rise
of \(D(a,b)\) does not refute an eventual bound
\(R_N\ll N^{5/3-\delta}\), since its constant and onset are not fixed here.
Third, a finite drop does not prove such a saving. Finally, the radius product
is not an arithmetic \(L^2\) estimate and does not reassemble the signed
prime-shell packet in Gate B.

\section{Conclusion}
TPC-270 adds the missing finite scale lens to the TPC-267--269 chain. The exact
sixth-power normalization and positive interval ratios expose strong finite
cross-scale variation, while matched profile controls remain in a separated
subunit band. The next source-level task is to state and prove an explicit
uniform radius upper bound with a declared power and clock. The actual V59
radius and phase, strict \(1/400\) payment, arithmetic \(L^2\), full Gate B, and
the twin-prime conclusion remain open.

\begin{thebibliography}{9}
\bibitem{hardywright} G.~H.~Hardy and E.~M.~Wright,
\emph{An Introduction to the Theory of Numbers}, 6th ed., Oxford University
Press, 2008.
\bibitem{montgomeryvaughan} H.~L.~Montgomery and R.~C.~Vaughan,
\emph{Multiplicative Number Theory I: Classical Theory}, Cambridge University
Press, 2007.
\bibitem{tpc269} L.~Wang, ``A growing-cutoff and convex-profile transfer audit
for the V59 residual,'' TPC-269 project release, 2026.
\end{thebibliography}
\end{document}
